\chapter{Input and Simulation Setup}
\label{app:input}

This appendix illustrates how the elements defined in
Chapter~\ref{chapter:proposed_model} are represented in the current LodusPop
implementation. A simulation setup comprises:
\begin{itemize}
    \item \textit{i)} JSON input files describing the environment, initial
    population, and routines (Sections~\ref{app:input_environment},
    \ref{app:input_population}, and~\ref{app:input_routine});
    \item \textit{ii)} an experiment configuration that selects the input
    files, simulation horizon, plugins, and plugin parameters
    (Section~\ref{app:experiment_configuration});
    \item \textit{iii)} Python plugins that extend the model vocabulary,
    routines, action registry, lifecycle services, or observations
    (Section~\ref{app:plugin}); and
    \item \textit{iv)} a Python simulator script that loads these components
    and advances the simulation (Section~\ref{app:simulator}).
\end{itemize}

The listings are deliberately small, but together they form one consistent
example. Separating structural inputs from experiment and plugin configuration
keeps the common definitions of points of interest, blobs, templates, routines,
and actions stable while making capacity, disruption, dependencies, and
adaptation explicit scenario choices.


\section{Environment Input Files}
\label{app:input_environment}

The environment input represents the regions $R$ and points of interest $P$
defined in Section~\ref{sec:environment}. A region has a \textit{name}, a
reference coordinate in \textit{lng\_lat}, and a collection of
\textit{points\_of\_interest}. Each point of interest has a
\textit{unique\_name}, a semantic \textit{poi\_type} corresponding to
$\tau_i$, a coordinate, and an optional \textit{attributes} mapping
corresponding to $\bm{\alpha}_i$. Complete POI identifiers use the form
\textit{Region//unique\_name}; for example,
\textit{ExampleRegion//pharmacy\_0}.

Listing~\ref{lst:environment_description} contains a home, a school, a
pharmacy, and the infrastructure POIs used in the dependency example below.
The keys in \textit{attributes} have no universal behavior merely because they
appear in the environment file. For example, \textit{capacity} and
\textit{water\_level} acquire meaning only when a loaded plugin interprets
them. This distinction implements the separation between the POI attribute
mapping and domain policy established in Section~\ref{sec:plugins}.

\begin{lstlisting}[language={Python},
caption={Example of an environment input file with typed POIs and attributes.},
captionpos=b, label={lst:environment_description}, basicstyle=\scriptsize]
{
  "regions": [
    {
      "name": "ExampleRegion",
      "lng_lat": [-51.19757, -30.15647],
      "points_of_interest": [
        {
          "poi_type": "home",
          "unique_name": "home_0",
          "lng_lat": [-51.19805, -30.15633],
          "attributes": {"external_id": "H-001"}
        },
        {
          "poi_type": "school",
          "unique_name": "school_0",
          "lng_lat": [-51.19750, -30.15621],
          "attributes": {"capacity": 120}
        },
        {
          "poi_type": "pharmacy",
          "unique_name": "pharmacy_0",
          "lng_lat": [-51.19820, -30.15655],
          "attributes": {"capacity": 40, "water_level": 3.5}
        },
        {
          "poi_type": "power_source",
          "unique_name": "power_source_0",
          "lng_lat": [-51.19910, -30.15590]
        },
        {
          "poi_type": "water_source",
          "unique_name": "water_source_0",
          "lng_lat": [-51.19940, -30.15560]
        },
        {
          "poi_type": "water_source",
          "unique_name": "water_source_1",
          "lng_lat": [-51.19680, -30.15570]
        }
      ]
    }
  ]
}
\end{lstlisting}

The time-dependent direct disable causes $C_i(t)$ are runtime state rather than
static booleans in this file. A flood, failure, or scheduled-closure plugin adds
and removes its own named cause, so removing one cause cannot erase another.
Similarly, dependencies $D_i$ can be supplied through a dependency-data plugin.
Listing~\ref{lst:dependency_description} specifies that the pharmacy requires
the power source and at least one of the two water sources. It therefore
illustrates both $D_i^{\mathrm{all}}$ and a minimum-of group in
$D_i^{\mathrm{min}}$. Names containing \textit{//} are complete identifiers;
an unqualified name is valid only when it resolves to exactly one POI.

\begin{lstlisting}[language={Python},
caption={Example of all-of and minimum-of POI dependency rules.},
captionpos=b, label={lst:dependency_description}, basicstyle=\scriptsize]
{
  "dependency_rules": {
    "pharmacy_infrastructure": {
      "all_of": [
        "ExampleRegion//power_source_0"
      ],
      "min_of": [
        {
          "minimum_enabled": 1,
          "nodes": [
            "ExampleRegion//water_source_0",
            "ExampleRegion//water_source_1"
          ]
        }
      ]
    }
  },
  "node_dependencies": {
    "ExampleRegion//pharmacy_0": ["pharmacy_infrastructure"]
  }
}
\end{lstlisting}

The simulator recomputes effective availability after a direct state change
and propagates dependency effects until the states stabilize. Unavailability
does not itself move, delete, or queue a blob. The handler consuming an action
must explicitly decide whether to suppress, retain, redirect, or otherwise
transform the affected demand.


\section{Population Input Files}
\label{app:input_population}

The population input defines the initial blobs $B=(TP,\bm{Tc},\bm{Sc})$
described in Section~\ref{sec:population}. The
\textit{default\_traceable\_characteristics} mapping supplies a default value
for every declared traceable characteristic. A blob may override a default in
its own \textit{traceable\_characteristics} mapping. Plugins may subsequently
extend this schema through $\Delta\Sigma_{\Pi}$; a plugin-provided default must
be applied to the population before actions use the new characteristic.

The \textit{sampled\_characteristics\_categories} mapping declares the
categories available for every sampled characteristic. Within each blob, the
counts for each sampled characteristic must sum to $TP$. These are marginal
distributions: the input does not assert a joint distribution between, for
example, age and occupation.

The \textit{initial\_population} mapping associates complete POI identifiers
with their initial blob collections $\bm{B}_{P}$. POIs with no initial
population may be omitted. In Listing~\ref{lst:population_description}, blob
$B_i$ contains 100 unvaccinated individuals because it uses the default
traceable value. Blob $B_j$ contains 20 vaccinated children and overrides that
default. Both the age and occupation distributions of each blob independently
sum to its total population.

\begin{lstlisting}[language={Python},
caption={Example of a population input file.},
captionpos=b, label={lst:population_description}, basicstyle=\scriptsize]
{
  "default_traceable_characteristics": {
    "vaccination_status": "unvaccinated"
  },
  "sampled_characteristics_categories": {
    "age": ["children", "youngs", "adults", "elders"],
    "occupation": ["student", "worker", "other"]
  },
  "initial_population": {
    "ExampleRegion//home_0": [
      {
        "total_population": 100,
        "traceable_characteristics": {},
        "sampled_characteristics": {
          "age": {
            "children": 20,
            "youngs": 10,
            "adults": 50,
            "elders": 20
          },
          "occupation": {
            "student": 30,
            "worker": 50,
            "other": 20
          }
        }
      },
      {
        "total_population": 20,
        "traceable_characteristics": {
          "vaccination_status": "vaccinated"
        },
        "sampled_characteristics": {
          "age": {
            "children": 20,
            "youngs": 0,
            "adults": 0,
            "elders": 0
          },
          "occupation": {
            "student": 20,
            "worker": 0,
            "other": 0
          }
        }
      }
    ]
  }
}
\end{lstlisting}

All active traceable characteristics participate in blob compatibility. Thus,
blobs at the same POI can merge only when their traceable values agree; the
sampled category counts of compatible blobs are added during the merge pass.


\section{Routine Input Files}
\label{app:input_routine}

The routine input represents the ordered pairs $(t,A)$ in the local routines
$Ro_{P_i}$ and the global routine $GRo$ described in
Section~\ref{sec:routines}. Each action contains a type identifier $N$, a
\textit{population\_template} encoding the traceable and sampled predicates of
$T$, and a \textit{values} mapping encoding the parameters $\bm{V}$ required by
the registered handler. The current JSON interface stores the requested
quantity $T_Q$ as \textit{values.quantity}; conceptually, that value and the
two predicate mappings together determine the population selected by the
action. The optional \textit{name} is an annotation and has no execution
semantics.

Listing~\ref{lst:routine_description} first schedules a local
\textit{move\_population} action at $t=8$. It requests up to 20 unvaccinated
children who are students from the blobs in \textit{home\_0} and moves the
selected population to \textit{pharmacy\_0}. At $t=10$, the global routine
invokes the plugin-defined \textit{example\_action} once for the environment;
the plugin in Listing~\ref{lst:plugin_script} expands it into a movement back
to \textit{home\_0}. If multiple actions have the same cycle step, their file
order is significant because later actions observe all earlier state changes.

\begin{lstlisting}[language={Python},
caption={Example of ordered local and global routine inputs.},
captionpos=b, label={lst:routine_description}, basicstyle=\scriptsize]
{
  "global_routine": [
    {
      "cycle_step": [10],
      "execution_scope": "once",
      "action": {
        "name": "return_pharmacy_visitors",
        "type": "example_action",
        "population_template": {
          "traceable_characteristics": {
            "vaccination_status": "unvaccinated"
          },
          "sampled_characteristics": {}
        },
        "values": {
          "region": "ExampleRegion",
          "node": "pharmacy_0",
          "quantity": 20
        }
      }
    }
  ],
  "routines": {
    "ExampleRegion//home_0": [
      {
        "cycle_step": 8,
        "action": {
          "name": "send_students_to_pharmacy",
          "type": "move_population",
          "population_template": {
            "traceable_characteristics": {
              "vaccination_status": "unvaccinated"
            },
            "sampled_characteristics": {
              "age": ["children"],
              "occupation": ["student"]
            }
          },
          "values": {
            "origin_region": "ExampleRegion",
            "origin_node": "home_0",
            "destination_region": "ExampleRegion",
            "destination_node": "pharmacy_0",
            "quantity": 20
          }
        }
      }
    ]
  }
}
\end{lstlisting}

A global action with \textit{execution\_scope=once} is generated once whenever
its cycle-step condition is satisfied. The alternative
\textit{execution\_scope=per\_node} expands the action once for each matching
POI and supplies that POI's region, type, name, and identifier to the action
values. Local and global routines are followed by any plugin-contributed
end-of-step actions and then by a single merge pass.


\section{Experiment Configuration}
\label{app:experiment_configuration}

The experiment configuration connects the structural inputs to the selected
extension set $\mathbb{P}$. It names the environment, population, and optional
routine files; defines the cycle length and number of cycles; and supplies the
parameter mapping $\bm{\theta}_{\Pi}$ for each plugin. A file entry may also be
a list, allowing several environment or population sources to be combined.

Listing~\ref{lst:experiment_configuration} selects the preceding files,
activates the dependency input from
Listing~\ref{lst:dependency_description}, and supplies the default destination
and quantity used by the example plugin. Naming a plugin configuration does
not instantiate its Python class by itself; the simulator script loads the
corresponding implementation explicitly.

\begin{lstlisting}[language={Python},
caption={Example of an experiment configuration file.},
captionpos=b, label={lst:experiment_configuration}, basicstyle=\scriptsize]
{
  "envgraph_inputs_files": {
    "environment_file": "Example-Environment.json",
    "population_file": "Example-Population.json",
    "routine_file": "Example-Routine.json"
  },
  "simulation_parameters": {
    "total_cycles": 3,
    "cycle_length": 24,
    "random_seed": 0
  },
  "node_dependency_data_plugin": {
    "dependency_files": ["Example-Dependencies.json"]
  },
  "example_plugin": {
    "destination_node": "home_0",
    "default_quantity": 20
  }
}
\end{lstlisting}


\section{Plugins}
\label{app:plugin}

Plugins implement the extension contract $\Pi=(\bm{\theta}_{\Pi},
\Delta\Sigma_{\Pi},\mathcal{R}_{\Pi},\mathcal{F}_{\Pi},
\mathcal{H}_{\Pi})$ defined in Section~\ref{sec:plugins}. A plugin may use only
the components required for its role. Listing~\ref{lst:plugin_script} uses
four of them: it reads $\bm{\theta}_{\Pi}$ from the experiment configuration,
adds an \textit{example\_state} traceable characteristic through
$\Delta\Sigma_{\Pi}$, and registers the pair
$(\textit{example\_action},f_{\textit{example}})$ in
$\mathcal{F}_{\Pi}$. Its time-step method is an empty lifecycle hook in
$\mathcal{H}_{\Pi}$, and it contributes no routine through
$\mathcal{R}_{\Pi}$.

The handler is composite: instead of changing the population directly, it
returns an ordered list containing a \textit{move\_population} action. That
action retains the original population template and is recursively dispatched
through the common $A=(N,T,\bm{V})$ contract. Consequently, the loaded plugin
set must also provide the \textit{move\_population} handler.

\begin{lstlisting}[language={Python},
caption={Example of a configured composite action plugin.},
captionpos=b, label={lst:plugin_script}, basicstyle=\scriptsize]
from core.plugin import ActionPlugin
from core.routine import Action


class ExamplePlugin(ActionPlugin):

    def load_plugin(self, simulation):
        self.simulation = simulation
        self.graph = simulation.env_graph
        self.config = simulation.experiment_config.get(
            "example_plugin", {}
        )

        # Delta Sigma: extend every blob with a default value.
        simulation.original_block_template.add_traceable_characteristic(
            "example_state", "idle"
        )
        self.graph.add_blobs_traceable_property(
            "example_state", "idle"
        )

        # F: register a composite action handler.
        simulation.add_action_type_to_function(
            "example_action", self.consume_example_action, False
        )

    def consume_example_action(
        self, pop_template, values, cycle_step, simulation_step
    ):
        destination = self.config.get(
            "destination_node", "home_0"
        )
        quantity = values.get(
            "quantity", self.config.get("default_quantity", 20)
        )

        move_values = {
            "origin_region": values["region"],
            "origin_node": values["node"],
            "destination_region": values["region"],
            "destination_node": destination,
            "quantity": quantity
        }
        return [Action("move_population", pop_template, move_values)]

    def update_time_step(self, cycle_step, simulation_step):
        pass

    def unload_plugin(self):
        pass
\end{lstlisting}

Primitive handlers may instead modify $\mathcal{S}(s)$ directly and return no
further actions. Routine plugins can generate start-of-step or end-of-step
actions, data plugins can expose shared services, and logger plugins can
observe state without adding domain fields to the core model. Regardless of
role, plugin requirements and shared action identifiers must be compatible
within the loaded extension set.


\section{Simulator Script}
\label{app:simulator}

The simulator script parses the experiment and its input files, loads the
selected plugins, and advances the lifecycle. In
Listing~\ref{lst:simulator}, the dependency-data plugin is loaded before any
plugin can change POI availability, the primitive movement handler is loaded
before the composite example is used, and the simulation then advances for the
horizon defined in the experiment configuration.

Each call to \textit{update\_time\_step} advances the absolute step $s$ and
derives its cycle step $t$. The simulator updates the applicable step-boundary
hooks and consumes actions in start-of-step, ordered global and local routine,
and end-of-step phases before finally merging compatible blobs. Direct or
queued action invocation remains available for programmatic scenarios, but it
uses the same action representation and dispatcher.

\begin{lstlisting}[language={Python},
caption={Example of a simulator script using the revised lifecycle.},
captionpos=b, label={lst:simulator}, basicstyle=\scriptsize]
from plugins.data.node_dependency_data_plugin import (
    NodeDependencyDataPlugin
)
from plugins.time_actions.move_population_plugin import (
    MovePopulationPlugin
)
from util.data_parse import generate_lodus_simulation

from example_plugin import ExamplePlugin


# Parse experiments/ExampleExperiment.json and its input files.
simulation = generate_lodus_simulation("ExampleExperiment")

# Load the extension set required by the configured actions.
simulation.load_plugin(NodeDependencyDataPlugin())
simulation.load_plugin(MovePopulationPlugin())
simulation.load_plugin(ExamplePlugin())

# Advance all cycles. The simulation object manages t and s.
number_of_steps = (
    simulation.time_status.total_cycles
    * simulation.time_status.cycle_length
)
for _ in range(number_of_steps):
    simulation.update_time_step()
\end{lstlisting}
